\chapter{Introduction} \label{ch:introduction}

Fast Radio Bursts (FRBs) are bright radio transients with durations of the order of milliseconds
that originate at cosmological distances~\cite{Petroff_2019}. Since the discovery of the first burst
in archival pulsar survey data~\cite{Lorimer_2007}, the number of detections has grown rapidly. In
particular, the Canadian Hydrogen Intensity Mapping Experiment (CHIME) has increased the known
population by several orders of magnitude~\cite{Amiri2018}, and its second FRB catalog already
contains more than $3600$ unique sources~\cite{CHIMEFRB2026}.

Beyond the open questions regarding their nature, FRBs are emerging as a novel cosmological probe.
Their radio signal is dispersed by free electrons along the line of sight, so that the observed
dispersion measure traces the diffuse, ionized baryons, which are otherwise difficult to
observe~\cite{McQuinn_2014}. Localized FRBs have already been used to measure the cosmic baryon
density~\cite{Macquart_2020} and the Hubble constant~\cite{Hagstotz_2022}, while forecasts indicate
that larger samples could constrain primordial non-Gaussianity~\cite{Reischke_2021} and baryonic
feedback~\cite{Sharma_2025}.

The astrophysical origin of FRBs, however, remains unresolved. The detection of an FRB-like burst
from the Galactic magnetar SGR 1935+2154 established a direct link to magnetars~\cite{CHIMEFRB_2020},
whereas FRBs localized in old stellar environments point toward additional, delayed formation
channels~\cite{Wang_2020, Petroff_2022}. Associating individual bursts with their host galaxies
requires arcsecond localizations, which are available only for a small fraction of the
population~\cite{RafieiRavandi_2020}. A complementary, statistical approach exploits the fact that
FRBs at cosmological distances trace the large-scale structure of the universe~\cite{MasuiSigurdson_2015}.
Their angular clustering therefore encodes the bias of the host population, which in turn depends
on the progenitor channel~\cite{Shirasaki_2017}. Because FRB samples are sparse compared with galaxy
catalogs, their auto-correlation is strongly affected by shot noise. Cross-correlating FRBs with
dense galaxy samples mitigates this limitation, and a statistically significant FRB-galaxy
cross-correlation has already been detected in the first CHIME/FRB catalog~\cite{RafieiRavandi_2021}.

Photometric galaxy surveys are particularly suited for such cross-correlations, as they provide
dense galaxy samples that can be divided into tomographic redshift bins. The Kilo-Degree Survey
(KiDS) has recently completed its fifth and final data release (DR5), referred to as KiDS-Legacy,
which provides calibrated redshift distributions for six tomographic bins~\cite{Wright_2024, Wright_2025}.

In this thesis, a forecasting pipeline is developed to assess how well the angular cross-correlation
between FRBs and KiDS-Legacy galaxies can constrain the clustering bias of the FRB host population.
Two progenitor scenarios are considered: young magnetars, whose formation closely follows the cosmic
star formation history, and neutron stars formed through delayed channels, each described by a
redshift-dependent linear bias. Both are combined with a shallow FRB survey resembling CHIME and a
deep survey resembling the Square Kilometre Array (SKA)~\cite{Hashimoto_2020}. The angular auto- and
cross-power spectra are computed in the Limber approximation and propagated into a Fisher forecast.
This makes it possible to quantify the gain of the joint analysis over the FRB auto-correlation
alone, as well as the prospects of distinguishing between the two progenitor scenarios.

The structure of this thesis is as follows: \autoref{ch:theoretical_background} reviews the
cosmological background, the properties and possible progenitors of FRBs, and the statistical
description of the large-scale structure. \autoref{ch:methodology} describes the modeling of both
tracer populations and the computation of the angular power spectra. The Fisher forecast is
presented and discussed in \autoref{ch:results}, while \autoref{ch:conclusion} summarizes the
findings and gives an outlook on future, wider galaxy surveys.